% !PRINTED 194
\weberpage{194}
%<<CONT
nur verträglich, wenn $\gamma = 0$ ist. Ebenso können $\alpha$ und $\delta$ nur
zugleich verschwinden.

Ist nun $\beta$ und $\gamma = 0$, so müssen $\alpha^2 + \delta^2$ und $i(\alpha^2 - \delta^2)$
reell sein, also $\alpha^2$ und $\delta^2$ conjugirt imaginär, und wegen (2) ihr
Product $= 1$. Je nachdem also in (2) das obere oder untere
Zeichen gilt, müssen $\alpha, \pm\delta$ conjugirt imaginär sein. Ebenso ergiebt
 sich, dass, wenn $\alpha$ und $\delta = 0$ sind, $\beta$ und $\mp\gamma$ conjugirt
imaginär sein müssen.

Nehmen wir also jetzt an, dass keine der vier Zahlen $\alpha, \beta, \gamma, \delta$
verschwinde. Dann müssen, damit die Grössen
\[
\begin{aligned}
&\alpha\delta+\beta\gamma,\quad i(\alpha\gamma+\beta\delta),\quad -i(\alpha\beta+\gamma\delta)\\
&\alpha\delta-\beta\gamma,\quad \alpha\gamma-\beta\delta,\quad \alpha\beta-\gamma\delta
\end{aligned}
\]
reell sein können,
\[
\begin{aligned}
&\alpha\delta \text{ und } \beta\gamma \text{ reell,}\\
&\alpha\gamma \text{ conjugirt imaginär mit } -\beta\delta,\\
&\alpha\beta \quad \text{\textquotedbl{}}\quad \text{\textquotedbl{}}\quad \text{\textquotedbl{}}\quad -\gamma\delta
\end{aligned}
\]
sein, und danach ist auch
\[
\alpha^2 = \frac{\alpha\gamma \cdot \alpha\beta}{\beta\gamma} \text{ conjugirt imaginär mit } \delta^2 = \frac{\beta\delta \cdot \gamma\delta}{\beta\gamma}
\]
und
\[
\beta^2 = \frac{\alpha\beta \cdot \delta\beta}{\alpha\delta} \quad \text{\textquotedbl{}}\quad \text{\textquotedbl{}}\quad \text{\textquotedbl{}}\quad \gamma^2 = \frac{\gamma\delta \cdot \alpha\gamma}{\alpha\delta};
\]
also muss $\alpha$ conjugirt mit $\pm\delta$, $\beta$ conjugirt mit $\mp\gamma$ sein. Dass
beide Male das gleiche Zeichen stehe, widerspricht der Forderung,
dass $\alpha\gamma$ und $-\beta\delta$ conjugirt sein sollen. Da das Product
zweier conjugirt imaginärer Grössen stets positiv ist, so gelten
die oberen oder die unteren Zeichen, je nachdem in (2) das
obere oder das untere Zeichen steht, d. h. je nachdem $A$ eigentlich
 oder uneigentlich orthogonal ist. Wir kommen also zu
folgendem allgemeinen Resultate, in dem auch die speciellen
Fälle von verschwindenden $\alpha, \beta, \gamma, \delta$ mit enthalten sind:

\gesperrt{Die nothwendige und hinreichende Bedingung
dafür, dass die orthogonale Substitution (1), (2) reell
sei, besteht darin, dass $\alpha, \pm\delta$ und $\beta, \mp\gamma$ zwei conjugirt
imaginäre Paare bilden, worin die oberen Zeichen bei
den eigentlich, die unteren bei den uneigentlich
orthogonalen Substitutionen gelten.}
